% !TeX root = ../../main.tex
\chapter[面向动态精度需求的状态感知混合精度计算与协同加速方法]{面向动态精度需求的状态感知\\混合精度计算与协同加速方法}\label{chap:evocg}
\section{研究动机}
\subsection{高精度稀疏迭代计算的性能开销}
前两章通过低比特筛选与候选范围缩减降低检索成本，其质量主要由候选覆盖与检索召回衡量。本章进一步研究由 SpMV 驱动的 CG 迭代过程。在此场景中，矩阵--向量乘积参与残差、步长和搜索方向的递推，局部近似不仅改变当前结果，还会影响未来计算状态。因此，研究重点由预设或查询相关的计算简化拓展到基于运行状态的精度控制。

CG 通常采用 FP64 以获得可靠的数值行为，而现代 GPU 支持多种更低位宽的数据格式。减少 SpMV 的数值传输量与算术成本具有潜力，但不能直接以低精度峰值吞吐比例预测求解加速。第\ref{chap:analysis}章的实验表明，SpMV 平均占 CG 求解时间的67.1\%，其余向量运算、归约和启动开销同样影响完整执行时间。因此，需要在加速 SpMV 的同时优化迭代内运行机制。

\subsection{低精度计算的数值误差与收敛影响}
对第 $k$ 轮 SpMV，式\eqref{eq:perturbed}中的扰动 $E_kp_k$ 会进入残差更新，随后改变搜索方向。浮点 CG 的递推残差不一定与真实残差 $b-Ax_k$ 完全一致，低精度可能扩大两者的间隙，也可能改变共轭关系和所需迭代次数\cite{greenbaum1989,simoncini2003,eshof2004}。

因此，本章从完整迭代轨迹评价数值质量，综合考察停止条件、迭代次数以及结束后以 FP64 重算的真实残差。低精度允许带来一定迭代增加，只要完整求解时间下降且质量满足要求；反之，单轮执行变快但总迭代过多，仍可能导致端到端减速。

\subsection{固定精度与静态调度方法的局限性}
固定 FP64 保留较保守的数值行为，却无法利用状态相关的低精度机会。统一低精度可能在数值敏感的问题上导致迭代增长或失败。仅依赖一个残差或元素幅值阈值的规则，虽然执行开销较低，但不一定能够识别不同矩阵、不同区域及不同迭代阶段的敏感性。

图\ref{fig:analysis-cg-traj}中的精度轨迹显示，有些矩阵在后期使用更多 BF16，而另一些问题在临近收敛时更保守，不能预设所有求解过程都应单调提高精度。已有混合精度 CG 系统探索了 Tile 执行、规则调度及离线选择\cite{yang2024mille,song2023callipepla,chen2025rl}，本章进一步研究与迭代状态交互的局部精度决策，并将其与可高效执行该决策的 GPU 基础结合。

\section{状态感知混合精度计算系统设计概述}
EvoCG 由 EvoCtrl、EvoSpMV 和 EvoRun 组成\cite{evocg2027}。EvoCtrl 从当前搜索方向提取局部状态，输出 FP64、FP32 或 BF16 动作；EvoSpMV 使用 Tile-CSC 布局和统一内核依据精度动作执行计算；EvoRun 将策略执行嵌入迭代流程，通过设备端推理、区间复用和融合执行控制额外开销。
\thesisfig[0.70]{evocg/overview.pdf}{EvoCG 的精度控制、稀疏执行与运行时模块}{fig:evo-overview}{%
图\ref{fig:evo-overview}中的反馈来自求解器本身：本轮精度影响 SpMV 及后续 CG 状态，新的搜索方向又用于下一次精度选择。默认设置在目标矩阵上离线训练策略，在线执行时使用固定策略参数进行推理，不在每轮求解时重新训练。跨矩阵复用和预条件迁移采用独立实验设置，分别考察策略在不同求解条件下的适用性。
}


与前两章通过候选筛选调整计算范围不同，EvoCG 保持矩阵的逻辑非零结构，以局部数值精度作为控制变量。精度决策同时影响当前计算成本和后续求解状态，因此需要在 CG 递推关系下联合优化数值行为与执行效率。

\section{基于强化学习的闭环精度控制方法}
\subsection{动态精度控制问题建模}
设第 $k$ 轮求解状态为 $z_k=(x_k,r_k,p_k)$，精度动作向量为 $a_k$。混合精度 SpMV 可表示为
\begin{equation}
 \widetilde q_k=(A+E_k(a_k))p_k,\qquad z_{k+1}=F(z_k,a_k).
 \label{eq:controlled-cg}
\end{equation}
其中 $F$ 由 CG 递推确定。动作既改变本轮执行成本，又通过 $E_k$ 改变下一状态。优化目标是最小化达到停止条件的求解时间：
\begin{equation}
 T(\{a_k\})=\sum_{k=0}^{N(\{a_k\})-1}C(a_k,z_k).
 \label{eq:control-objective}
\end{equation}
$N$ 依赖整个动作序列，故逐轮选择最快精度并不必然使总时间最小。

EvoCtrl 将这一交互过程按强化学习框架组织，以一次完整求解作为回合，以状态、动作及求解反馈构成训练样本。策略以搜索方向的局部表征作为完整求解状态的压缩观测，其信息量受特征选取方式限制。该选择的有效性由消融、收敛和迁移实验检验。

\subsection{求解状态表征与 Tile 级精度决策}
搜索方向直接参与 SpMV，同时通过 $p_{k+1}=r_{k+1}+\beta_kp_k$ 包含前序迭代的信息。EvoCtrl 将其划分为长度 $C=64$ 的连续片段，并对每个片段独立归一化：
\begin{equation}
 p_k^{(i)}=p_k[iC:(i+1)C],\qquad
 s_k^{(i)}=\frac{p_k^{(i)}}{\max(\norm{p_k^{(i)}}_2,\delta)},\quad\delta>0.
 \label{eq:evo-state}
\end{equation}
末尾不足64 个元素时补零。归一化降低全局尺度变化对网络输入的影响，但也压缩了绝对幅值信息，因此其适用性不能仅凭形式推导确定。

策略对每个片段输出三种动作之一，分别表示 FP64、FP32 和 BF16。一个片段对应矩阵的一组连续64 列，动作由该列块中的非空 Tile 共享。物理 Tile 还包括行分块维度，所以“Tile 级精度”在本实现中并不意味着同一列块中的每个行 Tile 都独立决策。明确这一粒度有助于理解数据布局与策略调用数量。

相邻列块可以选择不同精度，从而把较低成本的数值格式用于当前策略认为较不敏感的区域。FP64 路径始终保留，允许策略在困难状态下采取保守动作。实验中确有矩阵几乎全部使用 FP64，说明可用低精度比例受问题本身限制。

\subsection{强化学习策略训练与优化目标}
策略采用共享两层 ReLU 主干的轻量多层感知机，隐藏层大小为64 和64，策略头输出三类分布，价值头估计回报。训练使用近端策略优化（Proximal Policy Optimization，PPO）\cite{schulman2017}。其截断代理目标限制单次策略更新幅度，以减轻训练中动作分布变化过快的问题。

奖励同时考虑残差下降、执行成本和成功终止，原方案可表示为
\begin{equation}
 R_k=w_1\log\frac{\norm{r_{k-1}}_2}{\norm{r_k}_2}
 -w_2\frac{C(a_k)}{n_{\mathrm{tiles}}}
 +w_3\mathbf{1}\!\left[\frac{\norm{r_k}_2}{\norm{r_0}_2}\leq\tau\right].
 \label{eq:evo-reward}
\end{equation}
残差下降项反映求解进展，成本项反映低精度收益，终止奖励抑制只追求局部低成本而无法完成求解的策略。实验取 $w_1=10$、$w_2=0.2$、$w_3=30$，折扣因子0.99。数值实现需对零残差和非有限状态作相应处理；公式用于描述优化意图，并不提供训练收敛或求解收敛的理论证明。

\begin{algorithm}[htbp]\caption{EvoCG 策略训练的求解器交互过程}\label{alg:evo-train}
\begin{algorithmic}[1]
\REQUIRE 对称正定矩阵 $A$，右端项 $b$，策略网络及容差 $\tau$
\ENSURE 训练后的精度策略
\STATE 使用 FP64 参考求解估计迭代预算
\FOR{训练回合}
\STATE 初始化 $x_0=0$、$r_0=b-Ax_0$、$p_0=r_0$
\WHILE{未达到停止条件且未超出训练迭代预算}
\STATE 根据式\eqref{eq:evo-state}构造片段状态并采样精度动作
\STATE 按动作执行混合精度 SpMV
\STATE 执行 CG 步长、解、残差和搜索方向更新
\STATE 计算式\eqref{eq:evo-reward}的奖励并保存交互样本
\STATE 样本达到更新间隔时执行 PPO 更新
\ENDWHILE
\ENDFOR
\STATE 固定策略参数，转换为设备端推理表示
\end{algorithmic}\end{algorithm}

原方案在收集1024 个交互样本后更新，学习率为 $10^{-4}$，每次更新10 个 epoch，小批量128，截断比例0.2，GAE 参数0.95。训练回合上限为1000，训练阶段根据 FP64 参考迭代数设置预算。在线评估则使用统一的10000 次最大迭代限制，两种预算目的不同，不应混淆。

\section{面向混合精度 SpMV 的 GPU 执行优化}
\subsection{Tile-CSC 统一混合精度数据格式}
独立精度决策只有在执行系统能高效切换数值流时才有价值。若每种精度使用独立内核，可能需要反复分类、重排及启动；若在细粒度元素层面分支，又可能产生线程束控制流分歧。EvoSpMV 以固定 Tile 结构约束精度变化范围，使所有格式共享索引。
\clearpage
\thesisfig[1.00]{evocg/tilecsc.pdf}{Tile-CSC 与线程束级混合精度执行的协同设计}{fig:evo-tile}{%
矩阵划分为 $R\times C$ 的 Tile，其中 $C=64$，$R$ 在预处理时选择。Tile 之间按行块采用类似 CSR 的组织，\texttt{tileRowPtr}指出每个行块的非空 Tile 区间，\texttt{tileColIdx}记录列块编号。Tile 内部采用 CSC，列指针记录非零区间，行索引记录相对行块起点的偏移。图\ref{fig:evo-tile}展示这种外层行块、内层按列的组合。
}


每个非零位置保存 FP64、FP32 和 BF16 数值流，结构索引共享。精度选择主要改变读取的数值流，而不改变稀疏遍历。低精度转换使用列尺度，FP32 与 BF16 的指数位宽相同，使列缩放策略可以协调设计，但两者尾数精度及可表示值集合仍不同。为了向量化访问，列数据补齐至2 的倍数，并用哨兵标识填充零；补齐改变物理存储，不改变逻辑非零结构。

维护多份数值流增加存储容量，代价是避免在线反复转换和精度专用布局重构。该权衡更适合重复迭代和重复求解；对仅执行一次、存储极紧张的矩阵，额外准备和容量成本可能使其不具优势。

\subsection{Warp-centric 混合精度 SpMV Kernel}
EvoSpMV 为每个行块分配固定数量的线程束，各线程束跨步遍历该行块中的非空 Tile。一个线程束处理一个 Tile，32 个 lane 协同覆盖64 列，lane $i$负责第 $i$ 与第 $i+32$ 列。这样，输入向量、列尺度和列指针可获得较规则访问，而行索引与数值的不规则性限制在列内遍历。

同一 Tile 采用同一种精度，线程束内按统一路径选择数值流。不同 Tile 的精度可以不同，但不需要为每种精度单独启动内核。按列组织使同一查询分量和尺度能够被对应列的多个非零元素复用，减少重复读取。对于输出归约，系统仍需正确处理不同列对同一行的贡献，不能因为使用 CSC 而忽略写入合并成本。

这里的线程束规则化主要针对精度选择和任务遍历，并不意味着稀疏列长度差异完全消失。列内非零分布仍可能造成 lane 负载不均，因而格式设计还需与 Tile 大小和任务分配联合选择。

\subsection{精度异构执行与 GPU 资源优化}
由于同一列块共享精度，若外层也按列块连续安排任务，一个线程束组可能集中处理同类精度 Tile，导致组间成本差异。按行块组织时，每个组通常接触多个列块及其精度，有助于摊平精度引起的工作量变化。该机制将空间精度差异与已有稀疏负载分配结合。

预处理在 $R\in\{32,64,\ldots,16384\}$ 中选择行块高度，以改善行块间非空 Tile 数量和 Tile 内部非零分布。较小 $R$ 可以细化任务却增加 Tile 描述符，较大 $R$ 减少描述符却可能扩大局部负载差异。参数选择因此同时影响执行效率和存储占用，应以完整内核性能评估，而非只追求最少 Tile 数。

\section{面向迭代计算的低开销运行时优化}
\subsection{设备端精度决策}
将策略放在 CPU 上执行会引入状态回传和设备同步，容易抵消短迭代中的节省。EvoRun 将整个推理过程融合为设备端内核，权重离线量化为 FP16 并按 WMMA 使用方式重排，利用 Tensor Core 完成策略网络运算，以 FP32 累加。此处 Tensor Core 用于策略推理，不应将其峰值能力直接视为稀疏 SpMV 的实测吞吐。

融合推理减少中间特征写回和多个小内核启动，使状态构造、网络计算和动作输出能够紧密衔接。策略权重固定后在多次迭代中复用，在线成本由推理频率、列块数和实现效率决定。小矩阵或极短迭代中，控制开销占比仍可能较高。

\subsection{基于区间复用的精度控制调度}
相邻迭代的精度选择具有一定持续性，逐轮调用策略可能重复得到相同结果。EvoRun 利用离线轨迹选择复用间隔 $K^*$，在窗口内沿用动作。间隔候选为1、2、4、8 等，统计复用导致的升精度延迟比例，并选择该比例小于0.1 的最大间隔。

这一规则考虑了延迟升精度与延迟降精度的不同风险。前者可能在数值敏感状态下继续使用不足精度，后者主要损失潜在性能。因此，间隔并非仅以减少推理次数为目标。0.1 是轨迹分析阈值，不是对求解误差或失效概率的理论上界。在线状态与画像轨迹差异较大时，复用策略的有效性仍需重新验证。

\subsection{融合计算与 CUDA Graph 执行优化}
除 SpMV 外，CG 还包含内积、标量广播、残差及搜索方向更新。EvoRun 将相关操作组织为设备端合作式融合步骤，通过网格级同步完成必要归约和状态更新，减少独立内核启动及中间向量反复读写。SpMV 与融合步骤共同形成固定的执行基础，由 CUDA Graph 重放。
\thesisfig[1.00]{evocg/runtime.pdf}{设备端策略推理、区间复用与 CUDA Graph 运行机制}{fig:evo-runtime}{%
图\ref{fig:evo-runtime}体现控制路径与求解路径的协同。策略推理仍位于求解关键路径，其开销通过设备端融合与区间复用共同降低。融合步骤也需要满足合作式执行和同步条件。系统收益来源应区分统一稀疏内核、融合运行时与动态精度，下一节通过相同执行基础上的 FP64 对照进行隔离。
}


\section{实验评估与分析}
\subsection{实验环境与评估方法}
本章沿用 EvoCG 研究中的实验配置与评估结果\cite{evocg2027}。实验平台为 NVIDIA A100 40 GB，主机为 AMD EPYC 7K62 的8 核16 线程配置，内存93 GiB；系统为 Ubuntu 22.04，CUDA 12.2、GCC 11.4.0。比较对象包括基于 cuSPARSE/cuBLAS 的 FP64 CG、PETSc 3.24.5、Ginkgo 1.11.0，以及 Mille-FP64 和 Mille-Mixed。以上版本对应实验使用的软件环境。

所有方法使用 $x_0=0$、$b=A\mathbf{1}$，参考精确解为全1 向量。停止条件为递推相对残差 $\norm{r_k}_2/\norm{r_0}_2\leq10^{-10}$，最多10000 次迭代。结束后另用 FP64 重算真实残差。真实残差检查在实验的计时区间之外，因此在线时间不包含这一验证步骤。以相同右端项和初值评估保证了对比一致性，但未覆盖任意右端项分布。

47 个矩阵来自 SuiteSparse 集合\cite{davis2011}，覆盖八类科学应用，选择对称正定、维度不小于10000 且参考 CG 能在规定预算内完成的矩阵。集合未按条件数进行筛除，包含较困难问题；然而，矩阵选择仍限制结论范围，不能外推到不定、非对称或参考求解本身不收敛的系统。

默认混合精度结果采用目标矩阵专属的离线训练策略，在线时间不含训练。EvoCG-FP64 保留相同 Tile-CSC、内核融合和 CUDA Graph，仅将所有 Tile 设为 FP64。对 Mille 系列只在其发布实现成功完成的矩阵上汇总相应比较。因此，Mille 系列比较采用其成功求解的对应子集，其他基线在47 个矩阵上统计。

\subsection{端到端求解性能评估}
\thesisfig[0.84]{evocg/performance.pdf}{EvoCG 的端到端求解性能与跨系统加速比}{fig:evo-performance}{%
图\ref{fig:evo-performance}给出与各基线的完整求解性能。EvoCG 相对 cuSPARSE、PETSc、Ginkgo、Mille-FP64 及 Mille-Mixed 的几何平均加速比分别为2.52、3.06、2.48、1.93 和6.99，最大加速比分别为4.47、4.89、4.44、4.87 和22.49。表\ref{tab:evo-speedup}将两种统计量分别列出。
}

\begin{table}[htbp]\centering
\caption{EvoCG 相对各求解器的加速比}\label{tab:evo-speedup}
\begin{tabular}{lrr}\toprule
参考系统&几何平均加速比&最大加速比\\\midrule
cuSPARSE CG&2.52&4.47\\
PETSc&3.06&4.89\\
Ginkgo&2.48&4.44\\
Mille-FP64&1.93&4.87\\
Mille-Mixed&6.99&22.49\\\bottomrule
\end{tabular}
\end{table}

跨系统提升包含数据格式、内核映射、融合执行和精度策略的共同作用。Mille-FP64 同样具有融合执行优势，EvoCG 仍获得改进，说明针对精度异构的布局和调度具有额外价值。相对 Mille-Mixed 的较大提升还与其已发布实现的同步、类型转换和分支成本相关，因此，6.99 倍加速反映精度控制与执行实现的综合差异。

\subsection{动态精度控制与收敛性分析}
与相同执行基础的 EvoCG-FP64 比较，动态精度带来1.28 倍几何平均加速，平均迭代次数增加11.4\%。统计结果中平均56.9\%的 SpMV 计算使用 FP32 或 BF16。结束后的真实残差范数比值（Mixed/FP64）几何平均为1.0108，支持混合精度在测试集合上达到接近参考的数值质量。
\thesisfig[0.65]{evocg/ablation.pdf}{相同执行基础上的精度控制消融、迭代变化与真实残差}{fig:evo-ablation}{%
图\ref{fig:evo-ablation}为可读性展示21 个代表性矩阵，但上述汇总来自47 个矩阵。部分较适合低精度的问题只增加少量迭代，单轮节省几乎直接转化为总时间收益；另一些矩阵虽然明显增加迭代，仍因低精度比例较高获得较小加速。极端情况下，thermal2 相对相同执行基础上的 FP64 减速约8\%，gridgena 约减速2\%。这表明动态精度的单轮收益在部分矩阵上仍不足以抵消额外迭代和控制开销。
}


默认矩阵专属策略在47 个矩阵上均达到设定停止条件，未观察到发散。该结果验证了矩阵专属策略在所测矩阵、右端项和停止条件下的数值可行性。图\ref{fig:analysis-cg-traj}中的四个代表案例进一步说明不同数值敏感性：Dubcova2 迭代数不变，wathen100 由314 增至345，ted\_B\_unscaled 由41 增至45，bodyy4 由227 增至268。递推轨迹和真实残差应共同用于判断质量。

\subsection{关键技术消融实验}
为了分离策略与执行基础，实验用基于幅值阈值的精度规则替换 EvoCtrl，保留 Tile-CSC、混合精度内核与 EvoRun。对片段最大幅值 $m_i$，规则按目标阈值 $\epsilon=10^{-10}$ 分配 BF16、FP32 或 FP64：$m_i<0.1\epsilon$ 使用 BF16，$0.1\epsilon\leq m_i<\epsilon$ 使用 FP32，其余使用 FP64。

该规则在47 个矩阵中仅27 个于10000 次迭代内收敛，而矩阵专属 EvoCtrl 在全部47 个上完成求解。规则表现出保守与激进两种情况：保守时接近 FP64 迭代数但难以利用低精度，激进时可能增加大量迭代甚至失败。该对照表明，所采用的幅值阈值规则难以同时兼顾全部测试问题的低精度收益与收敛要求。

\begin{table}[htbp]\centering
\caption{相同执行基础上学习策略与幅值规则的代表性比较}\label{tab:evo-rule}
\begin{tabular}{lrrrr}\toprule
矩阵&EvoCG 迭代&规则迭代&EvoCG 时间（ms）&规则时间（ms）\\\midrule
finan512&58&57&2.68&4.06\\
bcsstm39&499&415&12.76&18.37\\
jnlbrng1&156&154&4.35&5.65\\
ted\_B&682&821&13.48&15.52\\
Dubcova2&193&10000&11.88&1571.47\\
Pres\_Poisson&3428&10000&131.05&560.07\\\bottomrule
\end{tabular}
\end{table}

表\ref{tab:evo-rule}中达到10000 次的规则配置未在上限内完成收敛，其耗时不能视为成功求解时间。bcsstm39 上学习策略比规则多84 次迭代，却将时间由18.37 ms 降至12.76 ms，说明迭代次数并非优化目标的唯一组成。相相同执行基础上的 FP64 对照和规则替换共同隔离了精度控制的作用，而跨系统结果反映完整协同系统的收益。

\subsection{预条件求解与策略泛化分析}
在 Jacobi 预条件 CG 中，实验复用各矩阵的非预条件训练策略，不重新训练。策略只控制 $Ap_k$，预条件应用及标量计算保持 FP64。47 个矩阵均达到停止条件，相对 cuSPARSE Jacobi-PCG 几何平均加速4.17 倍，相对相同执行基础上的 EvoCG-FP64 加速1.15 倍，平均迭代增加4.99\%。这支持策略在所测 Jacobi 变换下的复用，但不能推广为对任意预条件器均有效。

跨矩阵迁移采用另一实验：仅在306 阶、2018 个非零元素的 mesh2e1 上训练，将策略直接用于47 个未见矩阵。23 个矩阵获得性能提升，这23 个改善案例的几何平均加速为1.21 倍；另外15 个收敛但未加速，9 个不收敛。1.21 倍仅对应获得加速的23 个矩阵，完整迁移结果同时包含未加速和求解失败的情况。
\thesisfig[0.88]{evocg/transfer.pdf}{由小矩阵训练策略迁移到更大矩阵的代表性收敛案例}{fig:evo-transfer}{%
图\ref{fig:evo-transfer}展示五个成功收敛的代表案例，其效果存在差异。对分布接近的任务，直接复用可减少训练成本；迁移失败时，默认矩阵专属训练或经验证的 FP64 方案是必要回退。策略输入来自数值分布而非直接编码矩阵标识，有助于一定范围的迁移，但并未建立普适泛化保证。
}


\subsection{存储与运行时开销分析}
设备端策略推理在47 个矩阵上平均占在线时间1.09\%，中位数0.60\%；31 个矩阵低于1\%，38 个低于2\%。这一结果支持轻量网络、融合推理和区间复用能够使控制开销保持较低，但在完全不使用低精度的矩阵上，推理仍会产生净成本。
\thesisfig[0.75]{evocg/cost.pdf}{策略推理占比与 Tile-CSC 相对 CSR 的存储开销}{fig:evo-cost}{%
图\ref{fig:evo-cost}还表明，多份数值流与 Tile 补齐增加存储。实验对40 个非离群矩阵报告 Tile-CSC 总占用约为 CSR 的1.84 倍；其余7 个离群矩阵因非零元素分散于较多 Tile，结构与填充成本显著增大。1.84 倍反映40 个非离群矩阵的容量开销；对其余矩阵，Tile 分散程度与填充比例构成更显著的存储约束。
}


离线训练也是系统代价。mesh2e1 训练少于15 分钟，而目标矩阵专属训练通常需要数小时，适合重复求解同一或相近系统。若训练和预处理总成本为 $T_{\mathrm{off}}$，每次求解相对参考节省 $\Delta T>0$，则简单摊销估计要求重复次数满足 $N>T_{\mathrm{off}}/\Delta T$。该关系仅用于说明收支平衡，实际部署还需要考虑策略有效期、矩阵变化和验证开销。单次求解即便在线阶段加速，也未必具有总成本优势。

% !TeX root = ../../main.tex
\section{动态精度控制的数值与执行约束}
\subsection{真实残差与递推残差的关系}
有限精度迭代中，递推残差由前一轮残差减去近似乘积得到，而真实残差由当前近似解重新计算。定义残差间隙 $g_k=(b-Ax_k)-r_k$，忽略其他舍入项，仅考虑 SpMV 误差 $e_k=\widetilde q_k-Ap_k$，可得
\begin{equation}
 g_{k+1}=g_k+\alpha_ke_k.
\end{equation}
其来源是将 $x_{k+1}=x_k+\alpha_kp_k$和 $r_{k+1}=r_k-\alpha_k(Ap_k+e_k)$代入定义。由三角不等式得到
\begin{equation}
 \norm{g_N}\leq\norm{g_0}+\sum_{k=0}^{N-1}|\alpha_k|\norm{e_k}.
\end{equation}
这是忽略其他浮点误差的解释性界，不是完整浮点 CG 的误差分析；实际步长和向量更新也会引入舍入。

该关系说明，即使每轮误差都较小，也需要关注其沿迭代累积的影响。递推残差趋近零时，真实残差可能被间隙限制。结束后以 FP64 重算真实残差，可独立检验递推停止条件所对应的数值质量。若将这种检查加入在线系统，必须把额外 SpMV 成本计入部署时间。

\subsection{状态归一化的收益与信息损失}
归一化输入使网络更关注方向和局部分布，而不受搜索向量整体尺度变化支配。这样可以减轻训练中不同迭代数值范围的差异，并为跨尺度迁移提供可能。但不同绝对幅值的片段归一化后可能相同，其实际误差敏感性却不一定相同。

因此，归一化是一种特征选择，而非无损变换。可将片段范数、残差阶段或历史动作加入状态作为扩展，但这会增加推理输入和训练复杂度。比较不同状态设计时，应固定执行基础、网络预算和训练协议，避免把更大网络的优势误归因于特定数值特征。

\subsection{动作粒度与控制复杂度}
列块宽度 $C$决定策略调用数量，约为 $\lceil n/C\rceil$。较小 $C$提供更细的局部精度选择，却增加状态提取、动作存储和推理成本；较大 $C$降低控制开销，但可能让局部敏感区域迫使整个块使用高精度。Tile 行高 $R$又影响物理执行与元数据，两个参数的作用不同。

\begin{table}[htbp]\centering\caption{EvoCG 空间与时间粒度的作用}
\begin{tabularx}{\textwidth}{p{2.5cm}XX}\toprule
粒度参数&主要影响&需要联动观察的指标\\\midrule
列块宽度 $C$&动作覆盖范围和状态维数&精度比例、推理时间、收敛\\
行块高度 $R$&任务划分和 Tile 数量&负载分布、索引开销、补齐\\
每行块线程束数&并发 Tile 执行&占用率、同步与资源压力\\
复用间隔 $K^*$&策略进入关键路径的频率&升精度延迟、迭代数、控制成本\\\bottomrule
\end{tabularx}\end{table}

本实现固定64 列，与线程束协同遍历和网络结构相匹配。若改变 $C$，不仅需要重训练网络，还可能改变列指针长度、向量化访问和内核映射。因此，参数搜索应围绕完整系统进行，不能单独选择最好的网络输入长度后假设硬件成本不变。

\subsection{统一索引下的容量成本}
设逻辑非零数为 $z$，共享行索引每项 $b_i$字节。FP64、FP32、BF16 三份数值合计14 字节，单份 FP64 为8 字节。忽略 Tile 描述符与补齐，数值加索引的存储比近似为
\begin{equation}
 \rho_{\mathrm{base}}=\frac{14+b_i}{8+b_i}.
\end{equation}
例如以4 字节索引作说明，比例为1.5，而不是数值字段的1.75。实际 Tile-CSC 还包含列指针、Tile 行指针、描述符和填充，因此测得比例可能明显更高。

若补齐使物理非零槽位数为 $z_p\geq z$，这些数值流和索引成本应按 $z_p$计算。对于大量只含一个非零的列，补齐到2 会造成较高相对开销；若非零元素分散到更多 Tile，每 Tile 固定元数据又会放大总占用。实验离群矩阵与这一机制相符，但精确归因仍需逐项统计字节数。

\subsection{策略复用的保守性与失效条件}
精度动作在相邻迭代保持不变，才使区间复用具有收益。如果求解状态突然进入数值敏感阶段，长复用窗口可能错过及时升精度。原方案通过离线轨迹控制延迟升精度的比例，这是经验约束，而非对任意在线轨迹的防护。

可用于分析的指标包括每个列块的动作切换次数、升降精度比例、连续保持长度以及延迟切换对应的迭代增量。只报告平均推理占比不能说明复用策略是否合理。特别是极短求解任务，一次错误复用就可能占据很大比例，需要与长迭代问题分别分析。

部署时若右端项或初值改变，即使矩阵未变，状态轨迹也可能改变。此时是否继续使用相同 $K^*$需要验证。保守回退可缩短间隔或使用 FP64，但相应成本应进入系统评价，不能将失败检测与回退视为额外免费步骤。

\subsection{训练随机性与策略选择偏差}
强化学习训练包含动作采样和优化随机性，单次训练得到的策略不一定代表稳定效果。完整复现应记录随机种子、回合预算、检查点选择准则以及用于选择策略的验证负载。若仅报告若干训练中的最佳在线时间，可能引入选择偏差。

对矩阵专属策略，应说明训练和评价是否使用相同右端项；对跨矩阵策略，应保持测试矩阵未参与调参。本文所整理的原结果具有明确矩阵集合和默认训练设置，多随机种子下的稳定性和置信区间仍有待补充评估。后续重复实验应在固定协议下报告成功率与时间分布，而非只报告最优一次。

\subsection{失败样本的计分方法}
当某方法未在迭代上限内完成求解，其运行时间是失败耗时而非成功求解时间。把失败样本从全部平均中直接删除，会使保留子集偏向容易问题；把上限时间当作完成时间，又会混淆质量要求。因此，应先报告成功数量，再在明确的共同成功集合上比较时间。

\begin{table}[htbp]\centering\caption{求解实验中不同结果状态的报告方式}
\begin{tabularx}{\textwidth}{p{2.8cm}XX}\toprule
结果状态&质量信息&性能信息\\\midrule
按条件成功&递推残差及真实残差&成功求解时间、迭代次数\\
达到迭代上限&注明未满足停止条件&仅报告失败耗时，不计为成功时间\\
出现非有限数&记录首次发生位置&保留故障前耗时及状态\\
实现不支持&说明矩阵或格式约束&从该实现比较集合中明确排除\\
策略回退成功&分别记录回退前后质量&计入检测、回退与重算总时间\\\bottomrule
\end{tabularx}\end{table}

原论文对 Mille 系列的比较采用其发布实现成功完成的矩阵集合，跨矩阵迁移又包含9 个失败案例。二者都需要保留集合信息。本文将相同执行基础下的矩阵专属结果作为精度策略的主要证据，将迁移结果作为有边界的补充证据。

\subsection{预条件迁移与求解器结构变化}
Jacobi 预条件改变搜索方向的数值分布，但保持较简单的逐元素应用，因此提供一种受控的策略迁移测试。更复杂预条件器可能引入额外稀疏三角求解、不同精度误差和更高单轮成本，这时 SpMV 在总时间中的占比会改变。

若预条件使迭代次数显著减少，策略训练与推理成本的相对比例可能增加。若预条件计算成为主要开销，仅加速 SpMV 的收益又受 Amdahl 效应限制。因此，后续扩展需同时测量精度控制的数值效果和算子耗时构成，不能从 Jacobi 实验直接推断更复杂预条件器上的同等收益。

% !TeX root = ../../main.tex
\section{状态依赖误差的执行示例与验证方法}
\subsection{一个两步收敛的精确算术示例}
为说明 CG 状态之间的依赖，考虑二维对称正定系统
\begin{equation}
 A=\begin{bmatrix}1&0\\0&4\end{bmatrix},\qquad
 b=\begin{bmatrix}1\\1\end{bmatrix},\qquad x_0=0.
\end{equation}
下述计算使用精确分数运算，是解释算法的构造示例，用于分析 CG 的状态依赖关系。初始残差与搜索方向均为 $(1,1)^{\mathsf T}$。第一轮 $q_0=Ap_0=(1,4)^{\mathsf T}$，因此
\begin{equation}
 \alpha_0=\frac{r_0^{\mathsf T}r_0}{p_0^{\mathsf T}q_0}=\frac25,
 \quad x_1=\begin{bmatrix}2/5\\2/5\end{bmatrix},
 \quad r_1=\begin{bmatrix}3/5\\-3/5\end{bmatrix}.
\end{equation}
进一步得到 $\beta_0=9/25$，搜索方向为 $p_1=(24/25,-6/25)^{\mathsf T}$。第二轮步长为 $\alpha_1=5/8$，于是 $x_2=(1,1/4)^{\mathsf T}$，残差为零。该例展示精确算术下二维 CG 的有限步性质，实际浮点系统中的行为还受到舍入和实现细节影响。

\begin{table}[htbp]\centering\caption{二维示例中的精确 CG 状态}
\begin{tabular}{ccccc}\toprule
轮次&搜索方向&矩阵向量积&步长&更新后残差\\\midrule
0&$(1,1)$&$(1,4)$&$2/5$&$(3/5,-3/5)$\\
1&$(24/25,-6/25)$&$(24/25,-24/25)$&$5/8$&$(0,0)$\\\bottomrule
\end{tabular}\end{table}

这个过程中的后续方向由前一轮残差决定，而前一轮残差又依赖矩阵向量积和步长。因此，混合精度不是对独立 SpMV 任务分别作选择；同一轮的误差会改变下一轮输入，构成闭环控制需要考虑的状态演化。

\subsection{单次乘积扰动如何改变后续状态}
现在只在第一轮把乘积替换为 $\widetilde q_0=(1,9/2)^{\mathsf T}$，即人为加入 $e_0=(0,1/2)^{\mathsf T}$。该扰动是用于分析的构造值，与特定浮点格式的实际舍入误差无关。

新的步长为 $\widetilde\alpha_0=4/11$，更新解为 $(4/11,4/11)^{\mathsf T}$，递推残差为 $(7/11,-7/11)^{\mathsf T}$。但由原矩阵重新计算的真实残差是 $(7/11,-5/11)^{\mathsf T}$，二者间隙为
\begin{equation}
 g_1=\begin{bmatrix}0\\2/11\end{bmatrix}
 =\widetilde\alpha_0e_0.
\end{equation}
这与前述残差间隙关系一致。递推残差的范数并不等于真实残差范数，新的 $\beta$和搜索方向也与精确过程不同。即使下一轮恢复高精度，也是在已经改变的状态上继续执行。

该例有两个启示。第一，检查单轮乘积的相对误差不足以直接预测完整求解时间；它还通过步长、搜索方向和残差影响后续轮次。第二，终止时应明确检查的是递推量还是真实残差。学习策略可以利用状态相关性提高效率，但不能替代独立的数值质量验证。

\subsection{局部误差与矩阵结构的联系}
对混合精度乘积可写为 $\widetilde q=(A+\Delta A)p+\delta$，其中 $\Delta A$概括数值表示变化，$\delta$概括乘加及归约舍入。该分解仅是解释性记号，具体误差界需要根据实际内核和数据类型推导。相同位宽在不同 Tile 上产生的绝对扰动取决于矩阵值、搜索方向幅值、行内抵消和归约长度。

若某 Tile 中的搜索方向分量很小，该区域的乘积绝对贡献可能较低；但在存在强烈抵消的行中，即使每个局部贡献不大，其变化仍可能明显改变最后的行结果。因此，“当前值小”只能作为可能的精度容忍特征，不能作为普遍的安全降精度准则。状态特征与训练环境共同决定策略能否利用这些关系。

同样，矩阵条件数能够提示问题的数值敏感性，却不能单独决定某轮某 Tile 应使用何种精度。求解状态随迭代变化，局部向量分布也随之变化。本文采用搜索方向片段提供动态信息，但其状态表征没有包含所有可能影响收敛的因素，这也是迁移和鲁棒性需要实验验证的原因。

\subsection{精度比例的加权统计方法}
若一半 Tile 采用低精度，并不表示一半非零元素或一半计算时间采用低精度。Tile 之间的非零数和访问成本可能不同。可以分别定义按 Tile 数、按逻辑非零数以及按实际读取字节数加权的精度比例：三者对应决策分布、算术工作量和访存需求，回答的问题不同。

例如两个 Tile 分别包含100 和900 个非零元素，若仅第一个 Tile 采用 BF16，则按 Tile 数统计低精度比例为50\%，按非零数统计则为10\%。这里同样是构造示例。若存在补齐，物理槽位比例还会进一步变化。分析实验的精度比例时，应保持论文所采用的统计口径，不把它自动解释成等比例的带宽节省。

\begin{table}[htbp]\centering\caption{混合精度统计的不同加权口径}
\begin{tabularx}{\textwidth}{p{2.8cm}XX}\toprule
统计方式&反映的信息&解释限制\\\midrule
按动作或 Tile 数&策略选择的空间分布&不反映每个区域的大小\\
按逻辑非零数&参与运算的数据规模&忽略填充与元数据\\
按物理槽位数&实际数值流长度&仍未计入全部事务开销\\
按测得执行时间&关键路径上的成本构成&重叠执行时不宜直接相加\\\bottomrule
\end{tabularx}\end{table}

\subsection{学习策略与规则策略的公平比较}
比较学习策略和固定规则时，应保持它们可见的状态与使用的内核基础尽可能一致。若规则实现通过 CPU 决定动作，而学习策略完全在 GPU 运行，时间差中同时包含决策质量和执行位置差异。为了隔离决策能力，可补充同设备执行的规则；为了评价完整系统，则保留端到端实现并说明不同路径。

规则阈值也需要在独立验证集合上选择。如果学习策略经过大量训练，而规则仅采用一个未经校准的阈值，比较会混合调参预算差异。反过来，若规则使用测试结果逐矩阵挑选最佳参数，也会引入有利偏差。合理协议应提前规定搜索范围、训练或验证预算，以及最终测试时是否允许调整。

实验提供了规则基线的成功数量和时间对照，可用于说明特定规则在测试条件下的局限。该比较的结论限定于所采用的规则及其参数配置。学习策略的优势仍应通过相同执行基础上的比较、失败案例和训练成本共同评估，并与可验证的数值保护机制互补。

\subsection{运行时优化的收益归属}
统一内核、向量操作融合和 CUDA Graph 均可降低执行开销，即便所有矩阵值采用 FP64 也可能受益。因此，本章首先用 EvoCG-FP64 作为相同执行基础的参照，再比较动态精度。相对通用库的完整系统加速比用于展示系统效果，相同执行基础上的加速比用于分析精度控制的增量价值，两者分别刻画系统整体收益与精度控制的增量贡献。

进一步区分时，可以保持矩阵布局不变，分别启用策略推理、动作复用和低精度数值流。但消融配置可能改变执行重叠和缓存行为，各模块时间差不一定具有严格可加性。若将每个独立收益相乘，也可能重复计算相互依赖的改善。因此，评估采用列出完整配置组合及总时间，并用时间分解辅助解释。

CUDA Graph 还要求重复执行结构与数据依赖满足捕获和重放条件。动态动作若只改变设备数组中的精度标记，可以在固定图结构中被内核读取；若每轮动作决定完全不同的内核数量，则可能增加图更新成本。EvoCG 的统一执行路径使动态精度尽量表现为数据变化，这是算法粒度与运行时机制协同的具体体现。

\subsection{从矩阵专属策略到任务族复用}
同一矩阵配合不同右端项时，搜索方向轨迹仍可能改变；不同初值、停止阈值和预条件器也会改变状态分布。因而“矩阵相同”是策略复用的重要条件，却未必充分。部署前应明确训练配置与实际求解配置的差异，并观察成功率、迭代数和真实残差，而不只记录模型推理是否正常执行。

任务族复用可以降低训练摊销压力，但应把训练矩阵、验证矩阵和最终测试矩阵分开。若先观察所有测试矩阵再挑选一个迁移效果最好的源策略，其报告结果会包含选择偏差。已有跨矩阵实验揭示了可复用机会及失败边界，下一步需要更严格的任务族划分与多种子验证，而不能仅增加训练样本后假定泛化问题已经解决。

对求解失败的回退同样需要明确状态。可以从原初值使用 FP64 重新求解，也可以在当前近似解上重新计算真实残差并重启；两者具有不同成本和数值行为。这些回退方案的完整性能与数值验证有待后续研究。它们的共同要求是保留原问题和独立质量检查，使性能优化始终服务于相同的数值求解目标。

\section{本章小结}
本章将稀疏近似计算拓展至状态演化的 CG 迭代，提出以搜索方向局部特征驱动的闭环精度控制，并结合 Tile-CSC、统一 GPU 内核及低开销运行时实现。实验分别验证完整系统性能、相同执行基础上的精度收益、收敛质量、策略迁移和额外成本。结果支持在所测试条件下动态精度能够改善求解时间，同时明确矩阵专属训练、存储增长及跨矩阵迁移失败等边界。该研究与前两章共享计算需求驱动的协同设计思路，但采用独立的数值质量语义。
